\documentclass{article}

% Pacotes
\usepackage[utf8]{inputenc}
\usepackage[brazil]{babel}
\usepackage[letterpaper,top=2.5cm,bottom=2.5cm,left=3cm,right=3cm,marginparwidth=1.75cm]{geometry}
\usepackage{amsmath, graphicx, subcaption}
\usepackage[colorlinks=true, allcolors=blue]{hyperref}

\usepackage{listings}
\usepackage{xcolor}
\usepackage{tabularx}
\usepackage{array}
\usepackage{amssymb}
\usepackage{extarrows}
\usepackage{url}
\usepackage{hyperref}
\usepackage{booktabs}
\usepackage{float}
\usepackage{graphicx}
\usepackage{enumitem}


% Definições de estilo para Python
\lstdefinestyle{pythonstyle}{
      language=Python,
      basicstyle=\ttfamily\small,
      keywordstyle=\color{blue}\bfseries,
      commentstyle=\color{gray}\itshape,
      stringstyle=\color{orange},
      numbers=left,
      numberstyle=\tiny,
      stepnumber=1,
      numbersep=10pt,
      backgroundcolor=\color{white},
      frame=single,
      breaklines=true,
      tabsize=4,
      showstringspaces=false,
}


\begin{document}

% Capa personalizada
\begin{titlepage}
      \centering
      \includegraphics[width=4cm]{images/fig_logo_cefetmg.pdf}\par
      \vspace{0.5cm}

      {\LARGE Centro Federal de Educação Tecnológica de Minas Gerais \par}
      \vspace{0.2cm}

      {\Large Engenharia de Computação \\ Ciência de Dados \par}
      \vspace{4cm}

      {\Huge \textbf{Pré-Projeto} \par}
      \vspace{0.3cm}
      {\Large Preço do Café e da Soja $\times$ Dólar $\times$ Clima: \\
      uma análise de dados das commodities agrícolas mineiras \par}
      \vspace{0.8cm}

      {\LARGE Alvaro José Sem Pé \par}
      {\LARGE João Paulo da Cunha Faria \par}
      \vfill

      {\large Divinópolis - MG \par}
      {\large Setembro de 2026 \par}
\end{titlepage}

\tableofcontents
\newpage


\section{Introdução}

O Brasil é o maior produtor e exportador mundial de café e um dos maiores produtores de soja, e o
agronegócio responde por uma parcela significativa do Produto Interno Bruto (PIB) e da balança comercial
do país. Dentro desse cenário, Minas Gerais ocupa posição de destaque: o estado concentra cerca de metade
da produção nacional de café, distribuída principalmente nas regiões do Sul de Minas, do Cerrado Mineiro e
das Matas de Minas \cite{conab}. O tema, portanto, está diretamente ligado à realidade econômica da região
em que o trabalho é desenvolvido.

Commodities agrícolas como o café e a soja são negociadas no mercado internacional em dólar. Isso faz com
que o preço recebido pelo produtor brasileiro, em reais, dependa não apenas da oferta e da demanda
mundiais, mas também da taxa de câmbio. Além disso, por serem produtos agrícolas, sua oferta é fortemente
influenciada pelo clima: secas prolongadas, excesso de chuva ou geadas nas regiões produtoras podem reduzir
a safra e provocar fortes altas de preço, como ocorreu após a geada de julho de 2021 que atingiu as
lavouras de café de Minas Gerais.

Os dados necessários para estudar essas relações são públicos, porém estão espalhados em diferentes fontes
e formatos: cotações publicadas em páginas web, séries históricas em planilhas, dados cambiais disponíveis
por API e dados meteorológicos em bases científicas. A Ciência de Dados, entendida como a área
interdisciplinar voltada para a extração de conhecimento, detecção de padrões e obtenção de
\textit{insights} a partir de dados, oferece as ferramentas adequadas para coletar, integrar, tratar,
visualizar e modelar essas informações.

Este pré-projeto descreve a proposta de trabalho final da disciplina de Ciência de Dados, que consiste em
construir uma base de dados integrada com os preços do café arábica e da soja, a cotação do dólar e
variáveis climáticas das principais regiões cafeeiras de Minas Gerais, no período de 2020 até o presente,
e analisar como esses fatores se relacionam ao longo do tempo.


\section{Descrição do Problema}

O preço de uma commodity agrícola é resultado de muitos fatores simultâneos. Dois deles são
particularmente relevantes para o produtor brasileiro:

\begin{itemize}
      \item \textbf{Câmbio:} como o café e a soja são cotados em dólar no mercado internacional, uma
            valorização do dólar tende a elevar o preço em reais. Não se sabe, porém, \textit{quanto} dessa
            variação é repassada ao preço interno, nem \textit{com que atraso} esse repasse acontece;
      \item \textbf{Clima:} a produtividade das lavouras depende de chuva e temperatura adequadas. Eventos
            extremos (seca, geada) reduzem a oferta, mas o efeito sobre o preço não é imediato: a chuva de
            um período pode afetar a florada e, consequentemente, a safra meses depois.
\end{itemize}

Soma-se a isso uma dificuldade prática: os dados estão dispersos. Cada fonte possui formato, frequência e
granularidade diferentes (cotações apenas em dias úteis, clima diário, séries históricas em planilhas,
páginas HTML), o que exige um esforço considerável de coleta, limpeza e integração antes que qualquer
análise possa ser feita.

A inclusão da soja permite ainda uma comparação interessante: enquanto o café arábica é um produto
fortemente dependente das condições climáticas de uma região relativamente concentrada (Minas Gerais), a
soja é produzida em diversas regiões do país e tende a ser mais sensível ao câmbio e ao mercado
internacional. Assim, o problema pode ser resumido nas seguintes perguntas de pesquisa:

\begin{enumerate}[label=\textbf{P\arabic*.}]
      \item Como os preços do café arábica e da soja evoluíram entre 2020 e 2026?
      \item Existe sazonalidade nos preços, associada aos períodos de safra e entressafra?
      \item Os preços acompanham a variação do dólar? Com quanto tempo de defasagem?
      \item Períodos de pouca chuva ou de temperaturas extremas nas regiões produtoras são seguidos de alta
            no preço do café?
      \item Qual foi o impacto da geada de julho de 2021 sobre o preço do café?
      \item Qual das duas commodities é mais sensível ao clima e qual é mais sensível ao câmbio?
\end{enumerate}


\section{Objetivos}

\subsection{Objetivo geral}

Analisar a relação entre o preço do café arábica e da soja, a taxa de câmbio (dólar) e as condições
climáticas das principais regiões cafeeiras de Minas Gerais, no período de 2020 até 2026, por meio de
técnicas de coleta, tratamento, visualização de dados e aprendizado de máquina.

\subsection{Objetivos específicos}

\begin{itemize}
      \item Desenvolver um \textit{web crawler} em Python que colete diariamente, de forma automatizada, as
            cotações de café e soja publicadas no portal Notícias Agrícolas;
      \item Obter as séries históricas de preços (CEPEA/ESALQ), de câmbio (Banco Central do Brasil) e de
            clima (NASA POWER/INMET) para o período de 2020 a 2026;
      \item Tratar e integrar as diferentes fontes em uma única base de dados organizada por data;
      \item Criar variáveis derivadas que capturem efeitos defasados, como chuva acumulada e dólar de dias
            anteriores;
      \item Produzir visualizações que respondam às perguntas de pesquisa e disponibilizá-las em um painel
            interativo;
      \item Aplicar modelos de regressão linear e \textit{Random Forest} para identificar quais fatores mais
            explicam a variação dos preços;
      \item Manter todo o código, dados e notebooks em um repositório público e documentado.
\end{itemize}


\section{Metodologia}

A metodologia segue as etapas clássicas de um projeto de Ciência de Dados: extração, tratamento,
visualização e modelagem. As subseções a seguir descrevem as ferramentas e técnicas previstas para cada
etapa. Ressalta-se que algumas bibliotecas e ferramentas poderão ser alteradas ao longo do desenvolvimento
do trabalho, de acordo com as necessidades encontradas.

\subsection{Ferramentas e ambiente de desenvolvimento}

\begin{itemize}
      \item \textbf{Linguagem:} Python 3, por sua ampla adoção em Ciência de Dados e pelo grande número de
            bibliotecas disponíveis;
      \item \textbf{Ambiente:} Jupyter Notebook/JupyterLab para a análise exploratória e documentação, e
            Google Colab como alternativa para execução em nuvem e compartilhamento entre os integrantes;
      \item \textbf{Versionamento:} Git e GitHub;
      \item \textbf{Automação:} GitHub Actions, para executar o \textit{crawler} diariamente sem depender de
            uma máquina local ligada.
\end{itemize}

\subsection{Extração de dados}

Serão utilizadas três técnicas de extração: \textit{web scraping}, consumo de APIs e download de bases
históricas. A Tabela~\ref{tab:fontes} resume as fontes de dados previstas.

\begin{table}[H]
      \centering
      \small
      \caption{Fontes de dados, técnicas de extração e formatos.}
      \label{tab:fontes}
      \begin{tabularx}{\textwidth}{>{\raggedright\arraybackslash}p{2.6cm}
                  >{\raggedright\arraybackslash}X
                  >{\raggedright\arraybackslash}p{2.8cm}
                  >{\raggedright\arraybackslash}p{1.6cm}
                  >{\raggedright\arraybackslash}p{1.7cm}}
            \toprule
            \textbf{Fonte}          & \textbf{Dado}                                         &
            \textbf{Técnica}        & \textbf{Formato}                                      & \textbf{Frequência} \\
            \midrule
            Notícias Agrícolas      & Cotações diárias de café arábica e soja               &
            \textit{Web scraping}   & HTML                                                  & Diária (dias úteis) \\
            CEPEA/ESALQ             & Indicador do café arábica e da soja (Paranaguá), R\$/saca &
            Download de séries históricas & XLS/CSV                                         & Diária (dias úteis) \\
            Banco Central (SGS/PTAX) & Cotação do dólar comercial (venda) e IPCA            &
            API REST                & JSON                                                  & Diária / mensal     \\
            NASA POWER              & Precipitação e temperatura (mín., méd., máx.) por coordenada &
            API REST                & JSON/CSV                                              & Diária              \\
            INMET                   & Dados das estações meteorológicas das regiões produtoras &
            Download de arquivos    & CSV                                                   & Horária/diária      \\
            \bottomrule
      \end{tabularx}
\end{table}

\paragraph{\textit{Web scraping}.} O \textit{crawler} será um \textit{scraper} \textbf{direcionado}, pois
coleta dados de um conjunto fixo de páginas de um único site, e \textbf{contínuo}, pois será executado
diariamente para atualizar a base. Será implementado com as bibliotecas \texttt{requests} e
\texttt{BeautifulSoup}; caso alguma tabela seja gerada dinamicamente por JavaScript, será utilizado o
\texttt{Selenium}. Como a coleta ocorre apenas a partir do início do projeto, ela gera poucos meses de
dados; por isso, o histórico desde 2020 será obtido do CEPEA, e os dados do \textit{crawler} serão
cruzados com o CEPEA para validação.

\paragraph{APIs.} A cotação do dólar será obtida pela API do Sistema Gerenciador de Séries Temporais (SGS)
do Banco Central, e os dados climáticos pela API da NASA POWER, consultada nas coordenadas de municípios
representativos das três regiões: Sul de Minas (ex.: Varginha), Cerrado Mineiro (ex.: Patrocínio) e
Matas de Minas (ex.: Manhuaçu). Os dados do INMET serão usados de forma complementar, para validação.

\paragraph{Ética e legalidade.} Serão seguidas as boas práticas de \textit{web scraping} ético: respeito
ao arquivo \texttt{robots.txt}, identificação por meio de um \textit{User-Agent} descritivo e intervalos
entre requisições (\texttt{time.sleep()}) para não sobrecarregar o servidor. Os dados coletados são
cotações de mercado e dados meteorológicos, ou seja, \textbf{não são dados pessoais}, de modo que a Lei
Geral de Proteção de Dados (Lei nº 13.709/2018) não se aplica \cite{lgpd}. Ainda assim, os termos de uso
das fontes serão respeitados e as fontes serão sempre citadas.

O Código~\ref{lst:crawler} ilustra a estrutura prevista para o \textit{crawler}.

\begin{lstlisting}[style=pythonstyle, caption={Esboço do crawler de cotações.}, label={lst:crawler}]
import time
from datetime import date

import pandas as pd
import requests
from bs4 import BeautifulSoup

HEADERS = {"User-Agent": "CEFET-MG CD - projeto academico (contato: email)"}
URLS = {
    "cafe": "https://www.noticiasagricolas.com.br/cotacoes/cafe",
    "soja": "https://www.noticiasagricolas.com.br/cotacoes/soja",
}

registros = []
for produto, url in URLS.items():
    resposta = requests.get(url, headers=HEADERS, timeout=30)
    soup = BeautifulSoup(resposta.text, "html.parser")
    tabela = soup.find("table")
    for linha in tabela.find_all("tr")[1:]:
        colunas = [c.get_text(strip=True) for c in linha.find_all("td")]
        registros.append({"data": date.today(), "produto": produto,
                          "praca": colunas[0], "preco": colunas[1]})
    time.sleep(5)  # evita sobrecarregar o servidor

pd.DataFrame(registros).to_csv("data/raw/cotacoes.csv", mode="a",
                               header=False, index=False)
\end{lstlisting}

\subsection{Natureza e tipos dos dados}

Quanto à estrutura, os dados de origem se classificam como:

\begin{itemize}
      \item \textbf{Semi-estruturados:} páginas HTML do Notícias Agrícolas e respostas JSON das APIs, que
            possuem estrutura implícita e precisam ser interpretadas para extração dos valores;
      \item \textbf{Estruturados:} planilhas do CEPEA e arquivos CSV do INMET, com colunas bem definidas.
\end{itemize}

Após o tratamento, todos os dados serão convertidos em uma base \textbf{estruturada}, na forma de
\textbf{séries temporais} indexadas por data. A Tabela~\ref{tab:variaveis} apresenta as principais
variáveis e seus tipos.

\begin{table}[H]
      \centering
      \small
      \caption{Principais variáveis da base integrada.}
      \label{tab:variaveis}
      \begin{tabularx}{\textwidth}{>{\raggedright\arraybackslash}p{3.8cm}
                  >{\raggedright\arraybackslash}X
                  >{\raggedright\arraybackslash}p{3.8cm}}
            \toprule
            \textbf{Variável}           & \textbf{Descrição}                               & \textbf{Tipo}                       \\
            \midrule
            \texttt{data}               & Data da observação                               & Temporal                            \\
            \texttt{preco\_cafe}        & Indicador café arábica (R\$/saca de 60 kg)       & Quantitativa contínua               \\
            \texttt{preco\_soja}        & Indicador soja Paranaguá (R\$/saca de 60 kg)     & Quantitativa contínua               \\
            \texttt{dolar}              & Cotação do dólar comercial (R\$/US\$)            & Quantitativa contínua               \\
            \texttt{chuva\_mm}          & Precipitação diária (mm)                         & Quantitativa contínua               \\
            \texttt{temp\_min/med/max}  & Temperaturas diárias ($^\circ$C)                 & Quantitativa contínua               \\
            \texttt{regiao}             & Sul de Minas, Cerrado Mineiro ou Matas de Minas  & Qualitativa nominal                 \\
            \texttt{periodo\_safra}     & Safra ou entressafra                             & Qualitativa nominal (derivada)      \\
            \texttt{mes}, \texttt{ano}  & Componentes da data                              & Qualitativa ordinal (derivada)      \\
            \bottomrule
      \end{tabularx}
\end{table}

\subsection{Tratamento de dados}

O tratamento será feito principalmente com as bibliotecas \texttt{pandas} e \texttt{numpy}, e
compreende as seguintes etapas:

\begin{enumerate}
      \item \textbf{Padronização:} conversão de datas para um único formato, de valores textuais
            (``R\$ 1.234,56'') para numéricos, e padronização de unidades;
      \item \textbf{Integração:} junção de todas as fontes em uma única tabela indexada por data. Como não há
            cotações em fins de semana e feriados, mas há dados climáticos todos os dias, o clima será
            agregado (soma da chuva, média da temperatura) para os dias sem cotação ou as cotações serão
            propagadas (\textit{forward fill}), conforme a análise;
      \item \textbf{Dados ausentes e valores atípicos:} identificação de lacunas e de \textit{outliers}
            (por exemplo, pelo método do intervalo interquartil -- IQR), verificando se representam erros
            ou eventos reais (como a alta após a geada de 2021);
      \item \textbf{Engenharia de atributos:} criação de variáveis que capturem efeitos não imediatos, como
            chuva acumulada nos últimos 30 e 90 dias, dólar defasado em 1, 5 e 20 dias úteis, variação
            percentual diária e mensal dos preços, médias móveis e média climática das três regiões;
      \item \textbf{Deflação (opcional):} correção dos preços pelo IPCA, para comparar valores reais ao
            longo do período.
\end{enumerate}

\subsection{Visualização de dados}

A visualização será a parte central da análise. Serão utilizadas as bibliotecas \texttt{matplotlib} e
\texttt{seaborn} para gráficos estáticos e \texttt{plotly} para gráficos interativos. Estão previstos:

\begin{itemize}
      \item Gráficos de linha da evolução dos preços do café e da soja (P1);
      \item \textit{Boxplots} dos preços por mês, para identificar sazonalidade de safra e entressafra (P2);
      \item Gráficos com eixo duplo preço $\times$ dólar e gráficos de correlação cruzada com diferentes
            defasagens (P3);
      \item Gráficos de chuva acumulada e temperatura por região sobrepostos ao preço do café (P4);
      \item Destaque do período da geada de julho de 2021 na série de preços (P5);
      \item Mapas de calor (\textit{heatmaps}) de correlação entre todas as variáveis, separados por
            commodity (P6).
\end{itemize}

Ao final, os principais gráficos serão reunidos em um painel interativo desenvolvido com
\texttt{Streamlit}, permitindo filtrar período, commodity e região.

\subsection{Técnicas de \textit{machine learning}}

O aprendizado de máquina será utilizado como \textbf{ferramenta de análise}, com o objetivo de
identificar quais fatores mais explicam a variação dos preços, e não de construir um modelo de previsão
perfeito. Serão utilizados, por meio da biblioteca \texttt{scikit-learn}:

\begin{itemize}
      \item \textbf{Regressão Linear:} modelo simples e interpretável, cujos coeficientes indicam a direção
            e a intensidade da relação entre o preço e cada variável (dólar, chuva, temperatura e suas
            defasagens);
      \item \textbf{\textit{Random Forest} (regressão):} modelo capaz de capturar relações não lineares, cuja
            medida de importância das variáveis (\textit{feature importance}) será comparada com os
            coeficientes da regressão linear.
\end{itemize}

Por se tratar de séries temporais, a divisão entre treino e teste respeitará a ordem cronológica
(\texttt{TimeSeriesSplit}), evitando que dados futuros sejam usados para explicar o passado. Os modelos
serão avaliados pelas métricas MAE (erro absoluto médio), RMSE (raiz do erro quadrático médio) e
coeficiente de determinação ($R^2$). Os modelos serão treinados separadamente para o café e para a soja,
permitindo comparar a importância do clima e do câmbio em cada commodity.


\section{Repositório}

Todo o projeto será mantido em um repositório público no GitHub, disponível em
\url{https://github.com/<usuario>/<repositorio>}. O repositório garantirá o versionamento do código, a
colaboração entre os integrantes e a reprodutibilidade das análises. A estrutura prevista é:

\begin{verbatim}
projeto-cafe-dolar-clima/
|-- .github/workflows/   # GitHub Actions: execucao diaria do crawler
|-- crawler/             # scripts de coleta (scraping e APIs)
|-- data/
|   |-- raw/             # dados brutos, como foram coletados
|   |-- processed/       # base tratada e integrada
|-- notebooks/           # analise exploratoria, graficos e modelos
|-- app/                 # painel Streamlit
|-- requirements.txt     # dependencias do projeto
|-- README.md            # descricao e instrucoes de execucao
\end{verbatim}

Os dados brutos serão mantidos separados dos dados tratados, de forma que todo o processo de tratamento
possa ser refeito a partir das fontes originais. O \texttt{README.md} documentará as fontes de dados, a
forma de executar cada etapa e os principais resultados.


\section{Cronograma e Divisão de Tarefas}

A Tabela~\ref{tab:cronograma} apresenta o cronograma previsto para o desenvolvimento do trabalho.

\begin{table}[H]
      \centering
      \small
      \caption{Cronograma previsto.}
      \label{tab:cronograma}
      \begin{tabularx}{\textwidth}{>{\raggedright\arraybackslash}X >{\raggedright\arraybackslash}p{4.5cm}}
            \toprule
            \textbf{Etapa}                                                        & \textbf{Período}        \\
            \midrule
            Pré-projeto e criação do repositório                                  & Setembro de 2026        \\
            Desenvolvimento do \textit{crawler} e automação com GitHub Actions    & Setembro/Outubro de 2026 \\
            Coleta das séries históricas (CEPEA, Banco Central, NASA POWER/INMET) & Outubro de 2026         \\
            Tratamento, integração e engenharia de atributos                      & Outubro de 2026         \\
            Análise exploratória e visualizações                                  & Outubro/Novembro de 2026 \\
            Modelos de \textit{machine learning}                                  & Novembro de 2026        \\
            Painel Streamlit, relatório final e apresentação                      & Novembro/Dezembro de 2026 \\
            \bottomrule
      \end{tabularx}
\end{table}

As tarefas serão divididas da seguinte forma:

\begin{itemize}
      \item \textbf{Alvaro José Sem Pé:} desenvolvimento do \textit{crawler} e coleta de dados (Notícias
            Agrícolas, CEPEA, Banco Central e dados climáticos);
      \item \textbf{João Paulo da Cunha Faria:} tratamento, limpeza e integração dos dados;
      \item \textbf{Ambos:} análise exploratória, visualizações, modelos de \textit{machine learning},
            painel, relatório e slides.
\end{itemize}


\section{Conclusão e Pretensões}

Este pré-projeto propõe a construção de uma base de dados integrada e de uma análise sobre como o câmbio e
o clima influenciam os preços do café arábica e da soja, com foco nas regiões produtoras de Minas Gerais.
O trabalho percorre todas as etapas de um projeto de Ciência de Dados abordadas na disciplina: extração
por \textit{web scraping} e APIs, tratamento de dados de diferentes naturezas, visualização e aplicação de
técnicas de \textit{machine learning}.

Pretende-se, ao final do trabalho:

\begin{itemize}
      \item Disponibilizar uma base de dados limpa e integrada, de 2020 até 2026, que atualmente não existe
            reunida em um único lugar;
      \item Responder às perguntas de pesquisa, quantificando a relação entre preço e dólar (e sua
            defasagem), a sazonalidade dos preços e o efeito de eventos climáticos como a geada de 2021;
      \item Comparar o comportamento do café e da soja quanto à sensibilidade ao clima e ao câmbio;
      \item Entregar um painel interativo e um repositório documentado e reprodutível.
\end{itemize}

Como limitações, reconhece-se que o \textit{crawler} coletará apenas alguns meses de dados até o fim do
semestre, o que é contornado pelo uso das séries históricas do CEPEA, e que as relações encontradas
indicarão \textbf{correlação}, não necessariamente \textbf{causalidade}, já que o preço das commodities
também depende de fatores não incluídos no estudo, como estoques mundiais, demanda internacional e
produção de outros países.


\begin{thebibliography}{99}

\bibitem{conab}
COMPANHIA NACIONAL DE ABASTECIMENTO (CONAB). \textbf{Acompanhamento da Safra Brasileira de Café}.
Brasília: Conab. Disponível em: \url{https://www.conab.gov.br/info-agro/safras/cafe}. Acesso em: Setembro
de 2026.

\bibitem{cepea}
CENTRO DE ESTUDOS AVANÇADOS EM ECONOMIA APLICADA (CEPEA/ESALQ-USP). \textbf{Indicadores de preços:
café e soja}. Disponível em: \url{https://www.cepea.org.br/}. Acesso em: Setembro de 2026.

\bibitem{noticiasagricolas}
NOTÍCIAS AGRÍCOLAS. \textbf{Cotações}. Disponível em:
\url{https://www.noticiasagricolas.com.br/cotacoes/}. Acesso em: Setembro de 2026.

\bibitem{bcb}
BANCO CENTRAL DO BRASIL. \textbf{Sistema Gerenciador de Séries Temporais (SGS)}. Disponível em:
\url{https://www3.bcb.gov.br/sgspub/}. Acesso em: Setembro de 2026.

\bibitem{nasapower}
NASA. \textbf{POWER -- Prediction Of Worldwide Energy Resources}. Disponível em:
\url{https://power.larc.nasa.gov/}. Acesso em: Setembro de 2026.

\bibitem{inmet}
INSTITUTO NACIONAL DE METEOROLOGIA (INMET). \textbf{Banco de Dados Meteorológicos}. Disponível em:
\url{https://portal.inmet.gov.br/}. Acesso em: Setembro de 2026.

\bibitem{lgpd}
BRASIL. \textbf{Lei nº 13.709, de 14 de agosto de 2018}. Lei Geral de Proteção de Dados Pessoais (LGPD).
Disponível em: \url{https://www.planalto.gov.br/ccivil_03/_ato2015-2018/2018/lei/l13709.htm}. Acesso em:
Setembro de 2026.

\bibitem{mitchell}
MITCHELL, Ryan. \textbf{Web Scraping com Python: coletando mais dados da web moderna}. 2. ed. São Paulo:
Novatec, 2019.

\bibitem{pandas}
THE PANDAS DEVELOPMENT TEAM. \textbf{pandas Documentation}. Disponível em:
\url{https://pandas.pydata.org/docs/}. Acesso em: Setembro de 2026.

\bibitem{bs4}
RICHARDSON, Leonard. \textbf{Beautiful Soup Documentation}. Disponível em:
\url{https://www.crummy.com/software/BeautifulSoup/bs4/doc/}. Acesso em: Setembro de 2026.

\bibitem{sklearn}
PEDREGOSA, F. et al. Scikit-learn: Machine Learning in Python. \textbf{Journal of Machine Learning
Research}, v. 12, p. 2825--2830, 2011.

\bibitem{matplotlib}
HUNTER, J. D. Matplotlib: A 2D Graphics Environment. \textbf{Computing in Science \& Engineering}, v. 9,
n. 3, p. 90--95, 2007.

\bibitem{seaborn}
WASKOM, M. L. seaborn: statistical data visualization. \textbf{Journal of Open Source Software}, v. 6,
n. 60, p. 3021, 2021.

\bibitem{plotly}
PLOTLY TECHNOLOGIES INC. \textbf{Plotly Python Open Source Graphing Library}. Disponível em:
\url{https://plotly.com/python/}. Acesso em: Setembro de 2026.

\bibitem{streamlit}
STREAMLIT INC. \textbf{Streamlit Documentation}. Disponível em: \url{https://docs.streamlit.io/}. Acesso
em: Setembro de 2026.

\end{thebibliography}

\end{document}
